% Material moved from Section 2 of the full version (../sections/s2_model.tex) for the standard-length article.
% Input in Section S1 (supp_model.tex). Text copied verbatim; only references are adapted (M- = main article; no prefix = Supplementary Material).

\subsection{Symbols and terms}\label{supp:std-model-symbols}

Table~\ref{tab:nomenclature} lists the symbols and terms used throughout; the parameters of the case study are in Table~\ref{M-tab:case}.

\begin{table}[!htbp]
  \centering
  \caption{Symbols and terms used throughout; the parameters of the case study are in Table~\ref{M-tab:case}.}
  \label{tab:nomenclature}
  \footnotesize
  \begin{tabular}{@{}>{\raggedright\arraybackslash}p{0.21\textwidth}>{\raggedright\arraybackslash}p{0.58\textwidth}>{\raggedright\arraybackslash}p{0.14\textwidth}@{}}
    \toprule
    Symbol or term & Meaning & Where\\
    \midrule
    $n$; $\Omega$; $T$ & spindle speed in rpm; spindle angular speed; tooth period $2\pi/(N_t\Omega)$ & Eqs.~\eqref{M-eq:phi}--\eqref{M-eq:chip}\\
    $a$; $\alim(g)$ & axial depth of cut; linear (ZOA) stability limit at the gain $g$, $\alim(0)$ without control & Eq.~\eqref{M-eq:char}\\
    $g$; $g_1$, $g_2$, $k_d$ & DVF gain; gains of the delayed law & Eqs.~\eqref{M-eq:dvf}, \eqref{M-eq:comb-law}\\
    Forced orbit $x_p(t)$ & $T$-periodic settled motion below $\alim$, with the clip inactive & Eq.~\eqref{M-eq:xp}\\
    $b(t)$; $\bmax$ & forced command $-g\dot x_p(t)$; its peak, sampled 256 times per tooth period & Eq.~\eqref{M-eq:xp}\\
    $\bar b$; $u_0$ & rigorous upper bound of $\max_t\lvert b(t)\rvert$; free voltage (residual authority) $\Vmax-\bar b$ in the estimates & Sec.~\ref{M-sec:cert-lmi}\\
    Authority & the command range $\lvert V\rvert\le\Vmax$, of which the forced orbit takes $b(t)$ & Sec.~\ref{M-sec:model-orbit}\\
    $\kappa$; $\achk$ & depth margin of the authority condition (0.9); checked depth & Eq.~\eqref{M-eq:achk}\\
    Step entry; face entry; feed ramp & full depth and immersion from $t=0$, starting from rest; tool advancing into a face normal to the feed; feed per tooth rising from zero to $f_z$ over ten revolutions & Sec.~\ref{M-sec:bistab-prot}, \ref{supp:std-bistab-prot}\\
    Kick $\Delta v$; critical kick $\Delta v^*$ & velocity step added to $\dot x$ of the settled cut; the smallest that ends in chatter in the simulator & Sec.~\ref{M-sec:bistab-prot}\\
    Unsafe band & depths below $\alim$ at which some kick of \SI{3}{\milli\metre\per\second} to \SI{3}{\metre\per\second} ends in chatter & Sec.~\ref{M-sec:bistab-prot}\\
    Bistability & coexistence of the forced orbit with a bounded chatter attractor & Sec.~\ref{M-sec:bistab-attr}\\
    D1 & chatter classifier over a window of 300 revolutions & Sec.~\ref{app:simulator}\\
    $\xi$; $E_r$; $k^\top$ & state, regeneration matrix and feedback row of the collocation model ($k$ there is not the modal stiffness) & Eq.~\eqref{M-eq:coll}\\
    $a_{\mathrm{circ}}$ & depth below which the forced orbit of the averaged model attracts every solution & Prop.~\ref{M-prop:circle}\\
    C2; C3 & certificates on the collocation model of the averaged and of the time-periodic model & Sec.~\ref{M-sec:cert-model}\\
    Certified size $\beta$ & along $e=e_2$ the certified kick; along $-\xi_p(0)$, $\beta\ge1$ certifies the step entry & Eq.~\eqref{M-eq:lmi-size}\\
    $\beta_{\mathrm{req}}$; $a_{\mathrm{rob}}$; $a_{\mathrm{cap}}$ & required certified kick; basin-aware depth, up to which every grid depth meets it; assumed depth cap & Eqs.~\eqref{M-eq:admissible}, \eqref{M-eq:arob}\\
    $\tau_s$ & segment length of C3 & Sec.~\ref{app:lmi}\\
    DVF\,+\,D & DVF with the delayed difference term ($k_d=0$) & Eq.~\eqref{M-eq:comb-law}\\
    \bottomrule
  \end{tabular}
\end{table}
% src: src/feasibility.py (KAPPA = 0.9, N_T = 256); src/c1_unsafe_zone.py (N_PHASE = 16, BAND_DV 3 mm/s-3 m/s, ENTRIES: face entry R/f_z = 160 teeth = 40 rev, feed ramp 10 rev); src/tdboundary.py (NREV = 300); src/c2_certificate.py (nodes_for); src/c3_periodic.py (M_C = 16 segments); src/params.py (V_max = F_max/K_a = 9 V)

\subsection{Actuator: units, added damping and limits}\label{supp:std-model-act}

The actuator transfer function $G_a(s)$ of Section~\ref{M-sec:model-loop} is in newtons per volt. While the clip is inactive the loop is linear, with the closed-loop receptance of Eq.~\eqref{M-eq:Gcl}; at $g=\SI{1500}{\volt\second\per\metre}$ the feedback raises the damping of the structural mode from 1.34 to \SI{15}{\percent}.
% src: src/params.py, src/structure.py; verification/zoa_sdm/actuator_mode.json (|Ga|/K_a = 1.004 at 129.3 Hz; structural pole at g = 1500: 15.07 %, against 1.34 % at g = 0)

A proof-mass actuator is limited in three ways: by the voltage or current that its amplifier can supply, by the force that its coil and magnet can produce, and by the stroke of its proof mass. The model contains only the first, as a clip on the command applied before the actuator dynamics. The rated force $F_{\max}=\SI{27}{\newton}$ thus enters only through the voltage $\Vmax=F_{\max}/K_a$ and is not a strict force limit: $\lvert G_a\rvert$ is within \SI{0.4}{\percent} of $K_a$ at the structural mode but peaks at $3.4K_a$ at \SI{8.6}{\hertz}, just above the suspension mode at \SI{8.4}{\hertz}, so transients with low-frequency content can exceed $F_{\max}$; in step entries at $g=1500$ the force reaches \SI{29.5}{\newton}, and a clip at \SI{8.1}{\volt} keeps it at \SI{26.6}{\newton}. The proof-mass stroke needed for a given force falls with the square of the frequency above the suspension mode, so a stroke limit would act mainly on the slow part of an entry transient; neither stroke nor current limits are modelled. The lower clip is compared with the nominal one in Section~\ref{supp:std-delayed-robust}.
% src: src/structure.py (Ga); |Ga|/K_a: 1.0040 at 129.3 Hz, peak 3.371 at 8.60 Hz (recomputed); README §4.6 (clip before the actuator); verification/paper/support_runs.json (clip: peak actuator force 29.52 N at 9 V and 26.63 N at 8.1 V, max_fact of the deepest stable step-entry run at 2500 rpm, g = 1500)

\subsection{Force model: angle convention, surface memory and helix}\label{supp:std-model-force}

The immersion angle $\phi_j$ of Eq.~\eqref{M-eq:phi} is measured clockwise from the feed-normal direction; in half-immersion down milling a tooth cuts while $\phi_j\in[\phi_{\mathrm{st}},\phi_{\mathrm{ex}}]=[\pi/2,\pi]$ (modulo $2\pi$), which is indicated by $\gamma_j(t,z)=1$ ($\gamma_j=0$ otherwise). In the chip thickness of Eq.~\eqref{M-eq:chip}, the surface has been left by the preceding tooth one tooth period $T$ earlier; Section~\ref{app:simulator} replaces $x(t-T)$ by the deepest earlier pass, which matters once the tool leaves the cut. Both $f_s$ and the coefficient $aH(t)$ of Eqs.~\eqref{M-eq:split} and~\eqref{M-eq:aHt} are $T$-periodic. Over one tooth period every axial slice contributes the same mean, only shifted in time by the helix, so the mean coefficient is exactly linear in the depth \citep{altintas1995analytical} [Eq.~\eqref{M-eq:aHmean}], with $\alpha_{yy}$ the directional factor of the feed-normal direction. As the axial tooth pitch $p_a=2\pi R/(N_t\tan\eta)$ is comparable with the depths of interest, the helix keeps $aH(t)$ within about $\pm\SI{12}{\percent}$ of its mean near the stability limit, and $aH$ is constant when $a$ is a multiple of $p_a$ \citep{zatarain2006helix,insperger2006helix}.
% src: src/zoa.py (alpha_matrix), src/tdsim.c (force), src/c3_periodic.py (aH); verification/zoa_sdm/README.md §1 (derivation, alpha_yy = -1.428399); README §12.2 (+-12 % near the limit), §14.4 (pitch 12.57 mm, Table 1)

\subsection{Computation of the linear limit}\label{supp:std-model-zoa}

At a fixed speed, Eq.~\eqref{M-eq:char} gives the depth $a(\omega)$ as a function of the chatter frequency, and $a_{\lim}(g)$ is its smallest real positive value on a scan up to \SI{600}{\hertz} in steps of at most \SI{0.05}{\hertz}, excluding the poles of $a(\omega)$ at the tooth-passing harmonics. It agrees with semi-discretisation of the same averaged model and is conservative by at most about \SI{1}{\percent} with respect to the time-periodic helical model (Section~\ref{app:simulator}).
% src: src/zoa.py (F_GRID: 1200 points over 0.5-60 Hz, i.e. 0.0497 Hz, and 0.02 Hz over 60-600 Hz; pole guard; no root below 60 Hz for g < ~2.4e4 V s/m, verification/zoa_sdm/actuator_mode.json); README §2-3; verification/zoa_sdm/current_grid_check.json (table_96: 1.96e-5); README §3 and Appendix A (time-periodic model: ZOA conservative by 0-1.04 %)

\subsection{Forced orbit: closed form, sampling and dependence on the depth}\label{supp:std-model-orbit}

Below $\alim$ and with the command inside the clip, the deviation decays and the motion converges to a $T$-periodic forced orbit $x_p(t)$. On any $T$-periodic motion the regenerative term of Eq.~\eqref{M-eq:split} vanishes, so $x_p$ is the response of the linear loop \eqref{M-eq:Gcl} to the static force alone, and it has a closed form. Let $\hat f_s(\theta)=\sum_m\hat f_m\,e^{imN_t\theta}$ be the static force per unit depth of the slice at $z=0$, as a function of the spindle angle $\theta=\Omega t$, with $\hat f_m$ its Fourier coefficients over one tooth period. Every other slice carries the same force, delayed in spindle angle by its helix lag $z\tan\eta/R$ [Eq.~\eqref{M-eq:phi}], so the integral over the depth is exact [Eq.~\eqref{M-eq:Fm}]. With $\omega_T=2\pi/T=N_t\Omega$, the orbit and its command are given by Eq.~\eqref{M-eq:xp}, evaluated with $N_h=40$ harmonics. The peak forced command $\bmax=\max_t\lvert b(t)\rvert$ is taken over 256 samples per tooth period; the certificates of Section~\ref{M-sec:certificates} use a rigorous upper bound $\bar b\ge\bmax$ instead (Section~\ref{app:lmi}). The factor in Eq.~\eqref{M-eq:Fm} vanishes when $a$ is a multiple of $p_a/m$, so at a fixed speed the forced command is not monotone in the depth. The orbit persists unchanged in the clipped loop as long as $\bmax<\Vmax$; $b(t)$ is then the command that the actuator supplies in settled cutting, and $\Vmax-\bmax$ is the authority left for any deviation from the orbit.
% src: src/feasibility.py (slice_force_harmonics, force_harmonics, forced_velocity: N_HARM = 40, N_T = 256); README §8.2 (helix cancellation, pitch 12.57 mm); src/c3_periodic.py (closed-form bound with the tail beyond 40 harmonics)

\subsection{Case study: parameters, damping ratio, published cuts and unmodelled effects}\label{supp:std-model-case}

Table~\ref{M-tab:case} lists the parameters, all taken from \citet{ozsoy2026robust} except those marked as derived or chosen here, and Table~\ref{tab:nomenclature} the symbols and terms used throughout. The flexible direction, not stated in the text of \citet{ozsoy2026robust}, is taken normal to the feed, along the actuator axis; this reproduces their uncontrolled stability chart with a median deviation of \SI{0.04}{\percent} over \SIrange{400}{3000}{\rpm}. The damping ratio is \SI{1.34}{\percent} rather than the \SI{1.45}{\percent} of their Table~1, because the receptance, pole and stability chart computed in that paper are fitted by the former, and Table~2 of \citet{ozsoy2025frequency} gives \SI{1.34}{\percent} for the same structure; with \SI{1.45}{\percent} the uncontrolled chart deviates by \SI{8}{\percent} (median), whereas the gain map of Section~\ref{M-sec:feas-pub} changes by a median factor of 1.00.
% src: README §2, §4.1, §5.1, §6; verification/consistency/a_damping.json (fig6a peak -111.10 dB, model -111.10 dB at 1.34 % and -111.79 dB at 1.45 %; fig6b pole 1.3405 %; fig9 fit 1.3394 %); results/paper_vs_ours.json (p05: 0.037 % at 1.34 %, 7.99 % at 1.45 %); results/gainmap_zeta145.npz vs gainmap_main.npz (median gain ratio 0.995)

The model was checked against the cutting tests of \citet{ozsoy2026robust}, and against a second data set on the same set-up in Section~\ref{M-sec:model-second}. At the gains used in the tests, the linear limit classifies 17 of the 19 distinct published cuts as measured. The two exceptions are controlled cuts that chattered below the predicted limit, at \SI{1300}{\rpm} and \SI{10}{\milli\metre} ($0.80\alim$) and at \SI{2800}{\rpm} and \SI{5.5}{\milli\metre} ($0.76\alim$); at the first, the frequency at which the model loses stability, \SI{149.0}{\hertz}, is close to the strongest measured peak, \SI{149.5}{\hertz}. The clipped model does not make them chatter either: their forced commands at the published gains, 3.5 and \SI{3.0}{\volt}, leave more than \SI{5}{\volt} free, step and face entries settle at both, and the critical kick is 25--\SI{27}{\milli\metre\per\second} over 16 kick phases at \SI{1300}{\rpm} (23--\SI{24}{\percent} of the cutting impulse of one tooth period), whereas no kick up to \SI{3}{\metre\per\second} causes chatter at \SI{2800}{\rpm}. Saturation-induced finite-amplitude instability (Section~\ref{M-sec:bistability}) thus cannot explain the second cut, and explains the first only if the test contained a perturbation of that size, which neither simulated entry produces; otherwise the model overestimates the controlled limit there by at least 25 and \SI{32}{\percent}.
Where chatter frequencies can be compared, model and measurement differ by 0.5--\SI{4}{\hertz}. The measured spectra also contain spindle harmonics that are not multiples of the tooth-passing frequency (the seventh is the strongest peak of the stable cut at \SI{1300}{\rpm} and \SI{8}{\milli\metre}), a sign of runout \citep{insperger2008chatter}, which the model omits. Like \citet{ozsoy2026robust}, the model neglects the beam modes of the set-up at 23 and \SI{47}{\hertz}, and it uses the exact velocity, which they estimate by integration and zero-phase filtering. The flute length and overhang of the tool, the mode shape over the depth, the proof mass and its stroke, the sensor and the loop latency are not reported in \citet{ozsoy2026robust} and are not modelled: the tool is rigid, the mode uniform over the depth and the loop, outside Sections~\ref{M-sec:ba-valid} and~\ref{M-sec:delayed}, without delay. The depths and kicks reported below are predictions of this model.
% src: results/cuts_main.json (21 marks, 19 distinct cuts; pred_paper equals measured at 17/19; 1300 rpm, 10 mm: g_paper 1157.98 V s/m, alim_paper 12.478 mm (10/12.478 = 0.80), fc_paper 149.03 Hz; 2800 rpm, 5.5 mm: g_paper 679.55 V s/m, alim_paper 7.264 mm (5.5/7.264 = 0.76)); verification/targets/fft_peaks.csv (Fig10 B: tallest peak 149.5 Hz; Fig10 A: tallest peak 151.6 Hz, spindle harmonic 7); README §5.6 (12B 131.9 vs 133.3 Hz, 12D 134.0 vs 138.1 Hz; 10B 149.0 vs 149.5 Hz), §4.6 (velocity estimated by integration and zero-phase filtering in the demand check of the base paper), §7 (beam modes 23 and 47 Hz neglected as in the base paper); verification/targets/targets.md (Table 1 of the base paper: beam modes 23 and 47 Hz; no flute length, overhang, mode shape, proof mass, stroke, sensor, sampling rate or latency recorded)
% src: verification/paper/revision_runs.json (closed_form.published_cuts: bmax_V 3.525 at 1300 rpm/10 mm, g 1157.98, and 3.000 at 2800 rpm/5.5 mm, g 679.55, i.e. 9 - bmax = 5.47 and 6.00 V; tooth_impulse_Ns 2.234, mass 20.30 kg; published_cuts_entries: step and face entries, 9 V clip, full surface memory (NREG = 0), chatter false at both cuts; published_cuts_kicks_mps: 1300 rpm 0.02505-0.02673 m/s over 16 phases, 2800 rpm inf at all 16 phases, i.e. no kick of DV_GRID up to DV_HI = 3 m/s (src/c1_unsafe_zone.py); impulse ratio 0.02505 x 20.30 / 2.234 = 0.228 to 0.02673 x 20.30 / 2.234 = 0.243; overestimate 12.478/10 = 1.248 and 7.264/5.5 = 1.321)

\subsection{The second data set: full statement of the comparison}\label{supp:std-model-second}

\citet{ozsoy2025frequency} report further tests on the same set-up; their Table~2, Eq.~(2) and Section~3.2 give the structure, actuator, tool and cutting constants of Table~\ref{M-tab:case}, including the damping ratio of \SI{1.34}{\percent} inferred above. Two groups of their tests were digitised, each twice and independently, with identical readings (Section~\ref{supp:oz25-data}): 39 controlled cuts with DVF at $g=\SI{253}{\volt\second\per\metre}$ reported stable, marginal or chattering (their Fig.~6), and 56 tests at 1700--\SI{2200}{\rpm} made after the workpiece had been clamped again (\SI{128.1}{\hertz}, \SI{1.48}{\percent}, \SI{1.13e7}{\newton\per\metre}), with the measured amplitude of the command or the mark that the actuator saturated (their Figs.~10 and 12). They form a second data set on the same rig, not a second machine.

The linear model reproduces their uncontrolled and controlled stability charts to 0.04 and \SI{0.14}{\milli\metre} (RMS, outside the steep flanks of the lobes) for the first structure and to 0.06 and \SI{0.07}{\milli\metre} for the second, whose gain, not stated, a fit of the chart puts at \SI{251}{\volt\second\per\metre}, within \SI{1}{\percent} of the \SI{253}{\volt\second\per\metre} of the first group, which is used below. Of the 39 controlled cuts, the model and their chart classify the same 26 as measured. Nine of the others are stable cuts at 700--\SI{1100}{\rpm} above the predicted limit, where the cutting speed is only 0.6--\SI{0.9}{\metre\per\second} and process damping, which the model omits, may raise the limit; two are stable cuts at 2800 and \SI{2900}{\rpm}, and two, marginal and chattering at \SI{2100}{\rpm}, lie below the limit.
% src: verification/ozsoy2025/setup2025.py (Table 2 of the 2025 article = src/params.py; RECLAMPED 128.1 Hz, 1.48 %, 1.13e7 N/m), digitise_notes.md, points_fig6.json, points_fig10.json; verification/ozsoy2025/fit_gain.py/.json (fig6: uncontrolled rms 0.040 mm, controlled_g253 rms 0.137 mm; fig10_uncontrolled rms 0.062 mm; fig10_gain_fit g_best 251, rms_at_253 0.074 mm); compare2025.py/.json and a recomputation with the digitised controlled curve of Fig. 6 (26 of 39 classified as measured by both, the same 13 mismatches: stable at 700 (4, 5, 7 mm), 800 (5), 900 (10, 12), 1000 (12, 15), 1100 (4), 2800 (4), 2900 (5); marginal 2100/5 and chatter 2100/6); cutting speed pi D n: 0.59 m/s at 700 rpm, 0.92 m/s at 1100 rpm

For the 36 unsaturated tests of the second group, the settled forced demand $K_ag\max\lvert\dot x_p\rvert$ (Section~\ref{M-sec:model-orbit}; \SI{27}{\newton} corresponds to the \SI{9}{\volt} clip) is 0.90 times the measured amplitude (median; interquartile range 0.83--1.02), and within \SI{2.5}{\newton} of it (median \SI{0.85}{\newton}) at the 16 tests of up to \SI{5}{\milli\metre} at 1700--\SI{2000}{\rpm} (Fig.~\ref{M-fig:oz25}b). At \SI{15}{\milli\metre} it is about twice the measured value: like the frequency-domain prediction of \citet{ozsoy2025frequency}, the model misses the measured minimum of the command near \SI{15}{\milli\metre}, while the force of a \SI{45}{\degree} helix is most nearly constant at the axial pitch, \SI{12.6}{\milli\metre}; an effective helix of \SI{42}{\degree} would remove much of this error, which points to the axial distribution of the force or of the mode (Section~\ref{supp:oz25-sat}). At 2050--\SI{2200}{\rpm} the measured command at 2--\SI{3}{\milli\metre} is 1.1 to 2.3 times the prediction. The settled cut is predicted to saturate at 4 of the 20 saturated tests; the command reaches the clip during the simulated face entry at 12 of them, and later at two more, through simulated chatter above the controlled limit.
% src: verification/ozsoy2025/forced_map.py (demand_components), sim2025.py/.json, compare2025.py/.json (fig10.amplitude_peak: n 36, median_ratio_pred_over_meas 0.903, ratio_q25_q75 0.832-1.021; abs_error_depth_le5_rpm_le2000_N: n 16, median 0.85, max 2.46; demand_over_measured.depth_15: median 2.05 (1.23-2.64); measured over predicted at 2050-2200 rpm, 2-3 mm: 1.09-2.34 (recomputed from rows); saturation_settled_closed_form both 4 of 20; saturation_any_time_face_entry both 14, pred_only 2); helix2025.py/.json (mean_abs_log 0.240 at 45 deg, 0.159 at 42 deg; pitch 12.57 mm at 45 deg)

Relative to the linear limits of the model (Fig.~\ref{M-fig:oz25}a), the measured outcomes fall as follows. All 37 tests below the uncontrolled limit settled, 10 of them saturated. Between the uncontrolled and the controlled limit, 5 of the 7 saturated tests chattered or were marginal, and none of the 6 unsaturated ones did (one-sided Fisher exact test, $p=0.016$; for chatter alone, 2 of 7 against 0 of 6, $p=0.27$); above the controlled limit the 3 saturated tests chattered, while 3 unsaturated ones at 1700--\SI{1800}{\rpm} and 20--\SI{22}{\milli\metre}, above the model's and the published controlled limit, were stable. The model does not reproduce the 5 losses of stability between the limits: it predicts their settled command at 12--\SI{19}{\newton}, except at \SI{1950}{\rpm} and \SI{20}{\milli\metre} (\SI{33}{\newton}), both simulated entries settle, and, with kicks at four phases, a kick of at most \SI{3}{\metre\per\second} ends in chatter at only two of the five tests, both marginal, with critical kicks of 0.14 and \SI{0.40}{\metre\per\second}, whose impulses $m\Delta v^*$, with the modal mass of the re-clamped structure ($k/\omega_n^2=\SI{17.4}{\kilogram}$), are 1.7 and 2.3 times the cutting impulse of one tooth period [Eq.~\eqref{M-eq:impulse}]; at the other three, among them the chatter tests at 2050 and \SI{2100}{\rpm}, no kick up to \SI{3}{\metre\per\second} at these phases makes the model chatter. There the measured command exceeds the prediction, as noted, and \citet{ozsoy2025frequency} report disagreement with their own frequency-domain model at 2050--\SI{2200}{\rpm}, which they relate to loss of contact of the actuator or other nonlinearities; for the first structure, their time-domain simulation, whose details are not given, places the onset of saturation at \SI{2100}{\rpm} at \SI{4.5}{\milli\metre} (their Fig.~6), where the forced demand of the present model is \SI{13.5}{\newton} (\SI{14.2}{\newton} for the re-clamped structure). The saturation tests are thus consistent with the mechanism of this paper: below the controlled limit, losses of stability occurred only above the uncontrolled limit and only in tests whose command saturated. As a cut that chatters also drives the command into the clip, and the article does not say when during a test the command saturated, the data cannot show that the saturation came first. They support the forced demand at depths of up to \SI{5}{\milli\metre} at 1700--\SI{2000}{\rpm}, but the model does not predict these losses of stability.
% src: verification/ozsoy2025/compare2025.json (fig10.by_position: below a_lim(0) n 37, sat 10, unstable 0; between n 13; above a_lim(g) n 6, chatter 3, sat 3; between_sat_vs_outcome 5/2 against 0/6, fisher_one_sided_p 0.0163; rows: 1700/20, 1750/20, 1800/22 stable above a_lim(g)); kicks2025.py/.json (13 points between the limits: finite at 2000/10 (138.7 mm/s, marginal) and 1950/20 (396.5 mm/s, marginal), inf at the other 11); onset2025.py/.json (first structure, squares of the article's Fig. 6; 2100 rpm: article onset 4.5 mm, demand there 13.5 N, no model onset below a_lim(g) = 7.78 mm); re-clamped structure at 2100 rpm / 4.5 mm: forced_demand 14.15 N (feasibility.forced_demand after setup2025.apply_structure(RECLAMPED)); article Sec. 4.2 (disparity at 2050-2200 rpm, 'actuator loss of contact effects or other nonlinearities'); kicks at phases 0, 4, 8, 12 of 16 (kicks2025.py PHASES; crit_by_phase_mmps finite at all four phases at 1950/20 and 2000/10, none at 2050/5, 2100/5, 2150/5); impulse ratio m dv*/(Fbar_y T) with the re-clamped m = 17.44 kg and Eq. (16): 2.419/1.452 = 1.67 at 2000 rpm / 10 mm, 6.916/2.979 = 2.32 at 1950 rpm / 20 mm; Fisher one-sided for chatter alone 462/1716 = 0.269; clip during the face entry at 12 of the 14 run_clipped Sat tests, at 2050/10 and 2100/10 only after full engagement (sim2025_fig10.json entry_peak_demand_N 17.3 and 13.5 N, run peak 337 and 353 N)
